% 张三慧第四版V1第10章《热力学第一定律》独立候选解。
% 原书PDF371--375；功A为气体对外作正功，Q为吸入热正。
% 图示路径须继续比对原书及配套辅导扫描；不得视作专家验收。
\section*{第10章\quad 热力学第一定律：独立解析}

\paragraph{习题 10.1}
理想气体 \(pV=nRT\)；
以 \(A\) 为基准点 \(T_A=300\,\mathrm K\)：
\[
 \boxed{T_B=300\frac{2(20)}{2(10)}
 =600\,\mathrm K},\quad
 \boxed{T_C=300\frac{1(40)}{2(10)}
 =600\,\mathrm K},\quad
 \boxed{T_D=300\frac{1(20)}{2(10)}
 =300\,\mathrm K}.
\]
图中 \(BC\) 双曲线是等温线
\(pV=40\,\mathrm{atm\,L}\)。
若按 \(A\to B\to C\to D\) 连接，
气体作功
\[
 A_{AB}=2\,\mathrm{atm}(20-10)\,\mathrm L
 =20\,\mathrm{atm\,L},
\]
\[
 A_{BC}=(40\,\mathrm{atm\,L})\ln(40/20)
 =40\ln2\,\mathrm{atm\,L},
\quad
 A_{CD}=1\,\mathrm{atm}(20-40)\,\mathrm L
 =-20\,\mathrm{atm\,L}.
\]
故
\[
 \boxed{A_{A\to D}=40(101.325)\ln2
 \approx2.809\times10^3\,\mathrm J}.
\]
各状态 \((V,T)\) 为
\((10,300),(20,600),(40,600),(20,300)\)，
第一和第三段等压，第二段等温。
\textbf{图10.22原箭头审核：}
图上部分箭头可能表示反向行程；
上述数值明确是指定
\(A\to B\to C\to D\) 的路径，
若按箭头要求逆路径必须改变功符号。

\paragraph{习题 10.2}
由第一定律约定
\(\Delta U=Q-A\)。
\(a\to b\) 沿 \(acb\)：
\[
 \boxed{U_b-U_a=560-356=204\,\mathrm J}.
\]
换沿 \(adb\) 时
\[
 \boxed{Q_{adb}=\Delta U+220
 =424\,\mathrm J}.
\]
沿 \(b\to a\)，
\(\Delta U=-204\,\mathrm J\)，
外界对系统作功 \(+282\,\mathrm J\)
即气体对外功 \(A=-282\,\mathrm J\)，
于是
\[
 \boxed{Q_{ba}=\Delta U+A
 =-204-282=-486\,\mathrm J}.
\]
系统\emph{向外放出} \(486\,\mathrm J\) 热量。
状态函数 \(\Delta U\) 不随路径变化。

\paragraph{习题 10.3}
\(64\,\mathrm g\) 氧气
\(n=64/32=2\,\mathrm{mol}\)。
常温双原子气体
\(C_{V,m}=5R/2,C_{p,m}=7R/2\)；
温度差 \(\Delta T=50\,\mathrm K\)。
两过程均有
\[
 \boxed{\Delta U=nC_{V,m}\Delta T
 \approx2079\,\mathrm J}.
\]
定容过程 \(A_V=0\)，
\(\boxed{Q_V=\Delta U\approx2079\,\mathrm J}\)。
定压过程
\[
 \boxed{A_p=nR\Delta T\approx831\,\mathrm J},
 \qquad
 \boxed{Q_p=nC_{p,m}\Delta T\approx2910\,\mathrm J}.
\]
这是同一始末温度的气体，
定容与定压的\emph{内能变化}相同。

\paragraph{习题 10.4}
氮气 \(n=10/28\,\mathrm{mol}\)，
等压吸热
\(Q=nC_{p,m}(T_2-300)\)。
\[
 \boxed{T_2=300+\frac{1000}
 {(10/28)(7R/2)}
 \approx396.2\,\mathrm K}.
\]
若材料显著激发振动自由度，
需要修改常温双原子热容模型。

\paragraph{习题 10.5}
等压双原子氢气
\(Q_p=n(7R/2)\Delta T\)，
故
\[
 \boxed{n=\frac{6.0\times10^4}{(7R/2)50}
 \approx41.2\,\mathrm{mol}}.
\]
定压体积功
\[
 \boxed{A=nR\Delta T=\frac27Q_p
 \approx1.71\times10^4\,\mathrm J}.
\]
内能变化
\[
 \boxed{\Delta U=n(5R/2)\Delta T
 =\frac57Q_p\approx4.29\times10^4\,\mathrm J}.
\]
定容升温相同50K，吸热
\(\boxed{Q_V=\Delta U\approx4.29\times10^4\,\mathrm J}\)。
压强数值不影响这种用摩尔定压热容反算气量。

\paragraph{习题 10.6}
记始末为下标 \(1,2\)，
\(A=\int p\,dV,\ \Delta U=nC_{V,m}(T_2-T_1)\)，
\(Q=\Delta U+A\)。
\[
\begin{array}{c|c|c|c}
\text{过程}&A\ (\text{对外功})&Q\ (\text{吸热})&
\Delta U\\\hline
V=\mathrm{const}&0&nC_V\Delta T&nC_V\Delta T\\
T=\mathrm{const}&nRT\ln(V_2/V_1)&A&0\\
Q=0\ \text{可逆绝热}&nC_V(T_1-T_2)&0&
nC_V(T_2-T_1)\\
p=\mathrm{const}&p(V_2-V_1)=nR\Delta T&
nC_p\Delta T&nC_V\Delta T
\end{array}
\]
其中 \(C_V=C_{V,m},C_p=C_{p,m}\)。
对可逆绝热还可写
\(A=(p_1V_1-p_2V_2)/(\gamma-1)\)。
\emph{绝热不可逆过程}不能无条件套
\(pV^\gamma=\mathrm{const}\)。

\paragraph{习题 10.7}
氧气 \(n=1\)，
\(T_1=300\,\mathrm K,V_1=.010\,\mathrm{m^3}\)，
\(V_2=.030\,\mathrm{m^3}\)。
准静态绝热满足
\(TV^{\gamma-1}=\mathrm{const}\)，
故
\[
 \boxed{T_2=300(1/3)^{0.4}
 \approx193.3\,\mathrm K}.
\]
初压
\(p_1=nRT_1/V_1\approx2.494\times10^5\,\mathrm{Pa}\)，
末压
\[
 \boxed{p_2=p_1(1/3)^{1.4}
 \approx5.36\times10^4\,\mathrm{Pa}}.
\]
对外功与内能变化
\[
 \boxed{A=nC_V(T_1-T_2)
 \approx2.22\times10^3\,\mathrm J},\quad
 \boxed{\Delta U=-A\approx-2.22\times10^3\,\mathrm J}.
\]

\paragraph{习题 10.8}
初态
\(p_1=2\,\mathrm{atm}\)，
\(V_1=40\,\mathrm L\)，\(n=3\)，
\(T_1=p_1V_1/(nR)\approx325.0\,\mathrm K\)。
先可逆绝热压缩到
\(V_2=20\,\mathrm L\)，
\[
 \boxed{p_{\max}=2\cdot2^{1.4}\,\mathrm{atm}
 \approx5.278\,\mathrm{atm}},
\]
\[
 \boxed{T_{\max}=T_1\,2^{0.4}
 \approx428.8\,\mathrm K}.
\]
压缩阶段
\(Q_{12}=0\)，
\(A_{12}=nC_V(T_1-T_2)
 \approx-6.475\,\mathrm{kJ}\)。
第二段在 \(T_2\) 等温从20L扩至40L，
\[
 \boxed{Q_{23}=A_{23}=nRT_2\ln2
 \approx7.414\,\mathrm{kJ}}.
\]
全程
\[
 \boxed{Q_{\rm net}\approx+7.414\,\mathrm{kJ}},
\quad
 \boxed{A_{\rm net}\approx+0.939\,\mathrm{kJ}},
\]
\[
 \boxed{\Delta U\approx+6.475\,\mathrm{kJ}}.
\]
在 \(p\)-\(V\) 图中先沿斜率较陡的
绝热线向左上移动，
再沿高温等温线向右下返回原体积；
最后压强高于初始压强。

\paragraph{习题 10.9}
\textbf{原图10.24几何：
气缸初总体积30L，气体占20L，
其余10L为真空。}
自由扩散时由20L充至30L，
绝热且无外界功，故
\(T_{\rm mid}=T_1\)。
氮气 \(n=35/28=1.25\,\mathrm{mol}\)，
\[
 T_1=\frac{p_1V_1}{nR}
 =\frac{(2\times101325)(.020)}{1.25R}
 \approx390.0\,\mathrm K.
\]
自由扩散后
\(p_{\rm mid}=p_1(20/30)
 \approx1.333\,\mathrm{atm}\)。
随后准静态绝热从
\(V_{\rm mid}=30\,\mathrm L\)
膨胀至50L，
\[
 \boxed{T_f=T_1(30/50)^{.4}
 \approx317.9\,\mathrm K},
\]
\[
 \boxed{p_f=p_{\rm mid}(30/50)^{1.4}
 \approx6.61\times10^4\,\mathrm{Pa}}.
\]
第一阶段 \(A=Q=\Delta U=0\)，
第二阶段
\[
 \boxed{A_{\rm all}=nC_V(T_1-T_f)
 \approx1.87\times10^3\,\mathrm J},\quad
 \boxed{\Delta U_{\rm all}=-A_{\rm all}}.
\]
如果把自由扩散误当作可逆绝热，
会给错误中间温度与末压。

\paragraph{习题 10.10}
理想气体声速
\(u=\sqrt{\gamma RT/M}\)。
标准态 \(T=273.15\,\mathrm K\)，
氧气 \(M=.032\,\mathrm{kg/mol}\)，
声速 \(u=317.2\,\mathrm{m/s}\)，
\[
 \boxed{\gamma=\frac{Mu^2}{RT}
 \approx1.418}.
\]
其接近双原子理想气体理论值
\(\gamma=7/5=1.4\)，
实验值差异含测量条件和热容温度依赖。

\paragraph{习题 10.11}
以一小块空气作绝热气团，
令大气满足静力平衡
\(dp/dh=-\rho g\)，
理想气体 \(\rho=Mp/(RT)\)；
可逆绝热关系
\[
 \frac{T}{T_0}
 =\left(\frac p{p_0}\right)^{(\gamma-1)/\gamma},
\]
且 \(C_{p,m}=\gamma R/(\gamma-1)\)。
用 \(dp=(C_{p,m}p/(RT))\,dT\)
与上述静力平衡得
\[
 C_{p,m}\frac{dT}{dh}=-Mg,
\quad
 T=T_0-\frac{Mgh}{C_{p,m}}.
\]
代回绝热状态方程，
\[
 \boxed{p=p_0\left(
 1-\frac{Mgh}{C_{p,m}T_0}
 \right)^{\gamma/(\gamma-1)}}.
\]
有效范围限定括号内为正，
且假设气团局部沿干绝热线、
大气静力平衡；有潜热时不适用同一温度递减率。

\paragraph{习题 10.12}
原题主教材PDF物理页372—373给出
炮筒直径 \(d=.80\,\mathrm m\)、长
\(L=4.0\,\mathrm m\)，
初始燃气容积 \(V_1=2.0\,\mathrm{m^3}\)，
初始\emph{绝对}压强
\(p_1=2.7\,\mathrm{atm}\)，
表演者质量 \(m=70\,\mathrm{kg}\)，
绝热指数 \(\gamma=1.4\)。
出口时燃气容积为
\[
 V_2=V_1+\pi(d/2)^2L
 =2+\pi(.4)^2(4)
 \approx4.01062\,\mathrm{m^3}.
\]
绝热关系给
\[
 p_2=p_1(V_1/V_2)^\gamma
 \approx1.0193\,\mathrm{atm}.
\]
气体对活塞所作的\emph{总膨胀功}
\[
 A_g=\frac{p_1V_1-p_2V_2}{\gamma-1}
 \approx3.3232\times10^5\,\mathrm J.
\]
\textbf{必须扣除大气压力作功。}
主教材说明仅 \(75\%\) 的燃气推动功
转为推动者前进的有效机械功，
按配套辅导书采用的效率作用位置
并在之后扣除大气压 \(p_0=1\,\mathrm{atm}\)
对活塞的反功：
\[
 K_{\rm exit}=0.75A_g-p_0(V_2-V_1)
 \approx2.4924\times10^5-2.0373\times10^5
 =4.5515\times10^4\,\mathrm J.
\]
因此按\emph{题面打印值}严格计算：
\[
 \boxed{v=\sqrt{2K_{\rm exit}/70}
 \approx36.1\,\mathrm{m/s}}.
\]
\textbf{配套辅导书数据冲突：}
S PDF物理页164第10.12题演算过程中
实际使用 \(p_1=2.6\,\mathrm{atm}\)、
\(p_2\approx.99\,\mathrm{atm}\)、
\(m=75\,\mathrm{kg}\)，
而主教材题面打印
\(p_1=2.7\,\mathrm{atm},m=70\,\mathrm{kg}\)。
辅导书由这一组不同的数值与粗舍入
得到 \(K\simeq3.18\times10^4\,\mathrm J\)、
\(v\approx29\,\mathrm{m/s}\)。
两者差异不是相对论、热容模型造成，
而是\emph{原题与配套书选用数字不一致}；
本候选保留原题结果36.1m/s并明确报告
辅导书29m/s，等待最终编辑审核。
如改成把效率作用于扣除大气功\emph{之后}
的净作功，还会得到另一个数；
不能不声明效率约定。

\paragraph{习题 10.13}
稳定流动节流，忽略外部轴功、
热交换与上下游宏观流速/高度差：
单位时间进流做流动功 \(p_1\dot V_1\)，
出流流动功 \(p_2\dot V_2\)。
第一定律的控制体形式
\[
 E_1+p_1V_1=E_2+p_2V_2.
\]
定义焓
\(H=E+pV\)，于是
\[
 \boxed{H_1=H_2}.
\]
理想气体若定比热则
\(H=nC_{p,m}T\)，节流温度不变；
实际气体存在焦耳—汤姆孙效应。

\paragraph{习题 10.14}
按题目给出的\emph{两次从同一初态单独加热}
理解，等量电功输入热量 \(Q\)。
定容实验：
\[
 Q=nC_V(T_1-T_0)
 =\frac{C_V}{R}(p_1-p_0)V_0.
\]
等压 \(p=p_0\) 实验：
\[
 Q=nC_p(T_2-T_0)
 =\frac{C_p}{R}p_0(V_1-V_0).
\]
令两式相等并求 \(C_p/C_V\)，
\[
 \boxed{\gamma=\frac{C_p}{C_V}
 =\frac{(p_1-p_0)V_0}{p_0(V_1-V_0)}}.
\]
若误理解为第二次加热\emph{接续}
在第一实验终态 \(p_1,T_1\) 后，
分母将出现 \(p_1\)，
与原题待证公式不符；
故须明确两次是从同初始态比较。


\paragraph{习题 10.15}
多方过程 \(pV^n=C\)。
指数 \(n=0\) 是等压，\(n=1\) 对理想气体是等温，
\(n=\gamma\) 是准静态绝热，\(n\to\infty\) 为等容极限。
按本项目约定 \(A=\int p\,dV\) 为气体对外正功；
\(n\ne1\) 时，
\[
 \boxed{A=\frac{p_2V_2-p_1V_1}{1-n}}.
\]
若原教材式子写成分母 \(n-1\)，
它对应\emph{外界对气体作功}的相反号，
不应混用。
由 \(TV^{n-1}=\mathrm{const}\)，
\(A=\nu R\Delta T/(1-n)\)。
将 \(\Delta U=\nu C_{V,m}\Delta T\)
代入 \(Q=\Delta U+A\)，
\[
 \boxed{C_m=C_{V,m}+\frac R{1-n}
 =C_{V,m}\frac{\gamma-n}{1-n}}.
\]
特殊值：
\(C_m(n=0)=C_{p,m}\)；
\(n=1\) 等温 \(\Delta T=0\)，
相应有限吸热不能用有限 \(C_m\Delta T\) 表示；
\(n=\gamma\Rightarrow C_m=0\)；
\(n\to\infty\Rightarrow C_m=C_{V,m}\)。

\paragraph{习题 10.16}
绝热容器总容积 \(40\,\mathrm L\)，
初始两室各为 \(V_0=20\,\mathrm L\)，
各1摩尔氮气。
\[
 T_0=\frac{p_0V_0}R
 =\frac{(1.013\times10^5)(.020)}R
 \approx243.7\,\mathrm K.
\]
A室绝热压缩，任意时刻
\(p_A=p_0(V_0/V_A)^\gamma\)，
机械平衡给 \(p_B=p_A\)，
而 \(V_A=V_{\rm all}-V_B\)。
B室的过程曲线因此
\[
 \boxed{T_B(V_B)=
 \frac{p_0V_B}R
 \left(\frac{V_0}{V_{\rm all}-V_B}\right)^\gamma}.
\]
当 A室体积变为10L，
B室变为30L，
\[
 \boxed{T_A=T_0\,2^{.4}\approx321.6\,\mathrm K},
 \quad
 \boxed{T_B=\frac{p_0\,2^{1.4}(.030)}R
 \approx964.8\,\mathrm K}.
\]
整个容器对外不作体积功，
外界输入热量只施加在B室，
最终两室总内能变化等于该热量：
\[
 \boxed{Q_B=C_{V,m}(T_A+T_B-2T_0)
 \approx1.66\times10^4\,\mathrm J}.
\]
不能只计算B室自己的内能增加，
因为B还对A做了压缩功。

\paragraph{习题 10.17}
理想Brayton循环中
1→2、3→4绝热，
2→3定压吸热，
4→1定压放热。
\[
 Q_h=\nu C_p(T_3-T_2),\quad
 |Q_c|=\nu C_p(T_4-T_1),
\]
故
\[
 \boxed{\eta=1-\frac{T_4-T_1}{T_3-T_2}}.
\]
压强比 \(r_p=p_{\max}/p_{\min}\)，
可逆绝热线给
\(T_2/T_1=T_3/T_4=r_p^{(\gamma-1)/\gamma}\)。
因此
\[
 \boxed{\eta=1-r_p^{-(\gamma-1)/\gamma}}.
\]
当 \(r_p=10,\gamma=1.4\)，
\[
 \boxed{\eta=1-10^{-2/7}\approx48.2\%}.
\]

\paragraph{习题 10.18}
取柴油循环压缩比
\(r=V_1/V_2>1\)，
等压加热截止比
\(\rho=V'_1/V_2>1\)。
绝热压缩：
\(T_b=T_a r^{\gamma-1}\)。
等压吸热：
\(T_c=\rho T_b\)。
绝热膨胀至 \(V_1\)：
\(T_d=T_a\rho^\gamma\)。
等压吸热
\(Q_h=\nu C_pT_a r^{\gamma-1}(\rho-1)\)，
等容放热
\(|Q_c|=\nu C_VT_a(\rho^\gamma-1)\)。
于是
\[
 \boxed{\eta=1-\frac{\rho^\gamma-1}
 {\gamma r^{\gamma-1}(\rho-1)}
 =1-\frac{(V'_1/V_2)^\gamma-1}
 {\gamma(V_1/V_2)^{\gamma-1}(V'_1/V_2-1)}}.
\]
该理想公式以定热容和可逆绝热为前提。

\paragraph{习题 10.19}
三库温度为 \(T>T_2>T_1\)。
在最高温库吸热 \(Q_T>0\)，
中温库吸热 \(Q_2>0\)，
低温库放热 \(Q_1>0\)。
题给 \(Q_1=Q_2\)。
可逆循环克劳修斯等式
\[
 \frac{Q_T}{T}+\frac{Q_2}{T_2}
 -\frac{Q_1}{T_1}=0
\]
推出
\(Q_T=Q_2T(1/T_1-1/T_2)\)。
因 \(W=Q_T+Q_2-Q_1=Q_T\)，
总外部吸热 \(Q_h=Q_T+Q_2\)，
\[
 \boxed{\eta=\frac{W}{Q_h}
 =\frac{T(T_2-T_1)}
 {T(T_2-T_1)+T_1T_2}}.
\]
不能把该三库循环等同于仅在
\(T,T_1\) 两热库工作的卡诺循环。

\paragraph{习题 10.20}
第一台卡诺热机从高温库吸热 \(Q_h\)，
做功 \(W_1=\eta_1Q_h\)，
排给第二台热量
\(Q_m=(1-\eta_1)Q_h\)。
第二台做功
\(W_2=\eta_2Q_m\)。
合机效率
\[
 \boxed{\eta=\frac{W_1+W_2}{Q_h}
 =\eta_1+(1-\eta_1)\eta_2}.
\]
若热库温度 \(T_h>T_m>T_c\)，
代入
\(\eta_1=1-T_m/T_h\)、
\(\eta_2=1-T_c/T_m\)，
得
\[
 \boxed{\eta=1-\frac{T_c}{T_h}}.
\]

\paragraph{习题 10.21}
表层海水 \(T_h=298.15\,\mathrm K\)，
深层海水 \(T_c=278.15\,\mathrm K\)，
理想卡诺效率
\[
 \boxed{\eta_C=1-T_c/T_h=20/298.15
 \approx0.0671\;(6.71\%)}.
\]
若机械功率为 \(1.00\,\mathrm{MW}\)，
\(\dot Q_h=\dot W/\eta_C
 \approx1.491\times10^7\,\mathrm W\)，
排废热速率
\[
 \boxed{\dot Q_c=\dot Q_h-\dot W
 \approx1.391\times10^7\,\mathrm W}.
\]
按第（3）问，
热机全部吸热来自表层水冷却25℃至5℃；
采用 \(c_p=4180\,\mathrm{J/(kg\,K)}\)、
\(\rho=1000\,\mathrm{kg/m^3}\)，
\[
 \boxed{\dot V=\frac{\dot Q_h}
 {\rho c_p(25-5)}
 \approx0.178\,\mathrm{m^3/s}}.
\]

\paragraph{习题 10.22}
冷冻室 \(T_c=263.15\,\mathrm K\)，
室温 \(T_h=288.15\,\mathrm K\)，
逆卡诺制冷系数
\[
 \varepsilon_C=\frac{Q_c}{W}
 =\frac{T_c}{T_h-T_c}
 =263.15/25=10.526.
\]
消耗 \(W=1000\,\mathrm J\) 电功，
\[
 \boxed{Q_c=1.053\times10^4\,\mathrm J}.
\]
真实冰箱性能低于此可逆理论值。

\paragraph{习题 10.23}
室外 \(T_h=305.15\,\mathrm K\)，
室内 \(T_c=294.15\,\mathrm K\)。
漏热率
\[
 \dot Q_c=3.8\times10^4\,
 \mathrm{kJ/h}
 =3.8\times10^7/3600\,
 \mathrm W
 \approx1.056\times10^4\,\mathrm W.
\]
由逆卡诺制冷系数
\(\varepsilon_C=T_c/(T_h-T_c)\)，
最低输入机械功率
\[
 \boxed{P_{\min}
 =\dot Q_c\frac{T_h-T_c}{T_c}
 \approx3.95\times10^2\,\mathrm W}.
\]


\paragraph{习题 10.24}
热机两热库温度
\(T_1=483.15\,\mathrm K\)（210℃）、
\(T_2=288.15\,\mathrm K\)（15℃），
卡诺热效率
\[
 \eta_{\rm eng}=1-\frac{T_2}{T_1}
 \approx0.4036.
\]
每公斤燃煤释放
\(q=3.34\times10^7\,\mathrm J\)，
输出功
\(W=q\eta_{\rm eng}\approx1.348\times10^7\,\mathrm J\)。
热泵将低温河水热量送入温度
\(T_3=333.15\,\mathrm K\)（60℃）的供暖水，
理想制热系数
\[
 \varepsilon_h=\frac{T_3}{T_3-T_2}
 =\frac{333.15}{45}\approx7.403.
\]
所以
\[
 \boxed{Q_{\rm heat}=\varepsilon_hW
 \approx9.98\times10^7\,\mathrm J},
 \qquad
 \boxed{Q_{\rm heat}/q\approx2.99}.
\]
额外热量来自河水，
不违反能量守恒。

\paragraph{习题 10.25}
在图10.31冷冻循环
\(d\to a\to b\to c\to d\) 中，
\(a\to b\) 是80℃定压凝结，
\(c\to d\) 是5℃定压蒸发，
另外两段绝热。
气体对外功为正，
\(\mathrm{kPa\,m^3=kJ}\)。
高温段：
\[
 A_{ab}=2305(.00946-.0682)
 =-135.396\,\mathrm{kJ},
\quad
 \Delta U_{ab}=1171-1969=-798\,\mathrm{kJ}.
\]
所以
\[
 \boxed{Q_{ab}=\Delta U_{ab}+A_{ab}
 =-933.396\,\mathrm{kJ}}.
\]
向环境排热
\(\boxed{Q_h=933.4\,\mathrm{kJ}}\)。
低温段：
\[
 A_{cd}=363(.4513-.2202)
 =83.889\,\mathrm{kJ},\quad
 \Delta U_{cd}=1641-1015=626\,\mathrm{kJ},
\]
因此
\[
 \boxed{Q_c=Q_{cd}=709.889\,\mathrm{kJ}}.
\]
一个完整循环内能恢复原值，
压缩机输入功
\[
 \boxed{W_{\rm in}=Q_h-Q_c
 =223.506\,\mathrm{kJ}}.
\]
制冷系数
\[
 \boxed{\varepsilon=\frac{Q_c}{W_{\rm in}}
 \approx3.176}.
\]
高、低温热库温度
\(T_h=353.15\,\mathrm K,\ T_c=278.15\,\mathrm K\)，
逆卡诺上界
\[
 \boxed{\varepsilon_C=\frac{T_c}{T_h-T_c}
 =\frac{278.15}{75}\approx3.709}.
\]
实际值低于可逆上界。
本题是制冷系数，不是热机热效率。

\paragraph{习题 10.26}
设高温 \(T_1>T_2\)（低温），
容积比 \(r=V_1/V_2>1\)。
高温等温压缩 \(1\to2\) 输入功
\(W_{12}=\nu RT_1\ln r\)；
低温等温膨胀 \(3\to4\) 输出功
\(W_{34}=\nu RT_2\ln r\)。
两段等容过程不作体积功，
净输入
\[
 \boxed{W_{\rm in}=\nu R(T_1-T_2)\ln r}.
\]
若配备\emph{理想再生器}，
\(2\to3\) 冷却释放的
\(\nu C_V(T_1-T_2)\) 全部储存，
在 \(4\to1\) 升温时返还，
不与外界冷热库交换该部分显热。
低温库只在等温膨胀段
供给 \(Q_c=\nu RT_2\ln r\)，故
\[
 \boxed{\varepsilon_{\rm reg}
 =\frac{Q_c}{W_{\rm in}}
 =\frac{T_2}{T_1-T_2}},
\]
与逆卡诺循环相同。

若\emph{没有再生器}，
题设的 \(2\to3\) 等容冷却热量
\(\nu C_V(T_1-T_2)\)
排进\emph{低温库}，
而 \(4\to1\) 等容升温热量
由\emph{高温库}供给，
低温库的净放热必须减去
\(2\to3\) 收到的热量：
\[
 Q_{c,\rm net}=
 \nu RT_2\ln r-\nu C_V(T_1-T_2).
\]
有效制冷系数
\[
 \boxed{\varepsilon_{\rm no\,reg}
 =\frac{T_2}{T_1-T_2}
 -\frac{C_V}{R\ln r}
 <\varepsilon_{\rm reg}}.
\]
若这一净吸热非正，
无再生的外库接法不能有效制冷。
